% =====================================================================
%  1.4 Rasjonale uttrykk — teori og oppgaver
% =====================================================================

\del{1}{Forkorte, gange og dele}

\begin{ordliste}
rasjonalt uttrykk & En brøk med polynom i teller og nevner, for eksempel $\frac{x+7}{x-4}$. & \mittsprak\\
forkorte & Dele teller og nevner på den samme faktoren. & \mittsprak\\
vilkår & Det som må gjelde for at uttrykket er definert, for eksempel $x \neq 2$. & \mittsprak\\
faktor & Et tall eller uttrykk som ganges med noe annet. & \mittsprak\\
ikke definert & Vi kan ikke dele på 0. Da har uttrykket ingen verdi. & \mittsprak\\
\end{ordliste}

\begin{mustkunne}
\begin{minipage}[t]{0.52\linewidth}
\textbf{Vi kan bare forkorte faktorer.}
\[ \frac{a\cdot c}{b\cdot c} = \frac{a}{b}\cdot\frac{c}{c} = \frac{a}{b}\cdot 1 = \frac{a}{b} \]
Summer kan ikke forkortes. Prøv med tall:
\[ \frac{2+3}{4+3} = \frac{5}{7} \quad\text{men}\quad \frac{2}{4} = \frac{1}{2} \]
\end{minipage}\hfill
\begin{minipage}[t]{0.44\linewidth}
\textbf{Metode}\par
\begin{enumerate}[leftmargin=1.4em,itemsep=1pt,topsep=2pt]
  \item Skriv vilkår: nevneren er ikke 0.
  \item Faktoriser teller og nevner.
  \item Forkort like faktorer.
  \item Gange: gang tellere og nevnere. Dele: snu brøken og gang. Telleren i brøken du deler på, blir da nevner. Den kan heller ikke være 0.
\end{enumerate}
\end{minipage}

\medskip
Vilkåret gjelder fra starten og \emph{fortsetter å gjelde} etter forkorting.
Husk $\frac{x-3}{3-x} = -1$, fordi $3-x = -(x-3)$.
\end{mustkunne}

\begin{eksempel}[1 — forkort og skriv vilkår]
Forkort $\dfrac{x^2-9}{x^2+3x}$.

\begin{stegtabell}
$x^2+3x = x(x+3) \neq 0$ & Vilkår: $x \neq 0$ og $x \neq -3$.\\
$= \dfrac{(x-3)(x+3)}{x(x+3)}$ & Konjugatsetningen i teller, felles faktor $x$ i nevner.\\
$= \dfrac{x-3}{x}$, \ $x \neq 0,\ x \neq -3$ & Forkort $x+3$. Vilkåret står fortsatt.\\
Kontroll med $x = 1$: & $\frac{1-9}{1+3} = -2$ og $\frac{1-3}{1} = -2$. Det stemmer.\\
\end{stegtabell}
\end{eksempel}

\begin{eksempel}[2 — gange]
Regn ut $\dfrac{x^2-4}{3x} \cdot \dfrac{6x^2}{x+2}$.

\begin{stegtabell}
Vilkår: $x \neq 0$ og $x \neq -2$ & Nevnerne $3x$ og $x+2$ kan ikke være 0.\\
$= \dfrac{(x-2)(x+2)}{3x} \cdot \dfrac{6x^2}{x+2}$ & Faktoriser telleren.\\
$= \dfrac{(x-2)\cdot 6x^2}{3x}$ & Forkort $x+2$.\\
$= 2x(x-2)$ & Forkort $\frac{6x^2}{3x} = 2x$.\\
Kontroll med $x = 3$: & $\frac{5}{9}\cdot\frac{54}{5} = 6$ og $2\cdot 3\cdot 1 = 6$. Det stemmer.\\
\end{stegtabell}
\end{eksempel}

\begin{vanligfeil}
\begin{itemize}[leftmargin=1.2em,itemsep=1pt,topsep=0pt]
  \item $\frac{x+2}{x} = 2$ er feil. $x$ er ikke en faktor i telleren. Du kan ikke forkorte ledd.
  \item $\frac{x+3}{x+5} = \frac{3}{5}$ er feil. Prøv $x = 1$: $\frac{4}{6} \neq \frac{3}{5}$.
  \item $\frac{x-3}{3-x} = 1$ er feil. Det riktige er $-1$, fordi $3-x = -(x-3)$.
  \item Å glemme vilkåret etter forkorting. $\frac{x-3}{x}$ er bare lik det opprinnelige uttrykket når $x \neq 0$ og $x \neq -3$.
\end{itemize}
\end{vanligfeil}

\nivaahode{1}

\oppg{1.1}{1}{\uh}
Forkort. Skriv vilkår.
\begin{tasks}(4)
  \task $\dfrac{6x^2}{9x}$
  \task $\dfrac{3x+6}{x+2}$
  \task $\dfrac{x^2-5x}{x}$
  \task $\dfrac{x-4}{x^2-4x}$
\end{tasks}

\oppg{1.2}{1}{\uh}
Vi skal forkorte $\dfrac{x^2-25}{x^2-5x}$.
\begin{tasks}(2)
  \task Faktoriser telleren og nevneren.
  \task Forkort og skriv vilkår.
\end{tasks}

\oppg{1.3}{1}{\uh}
Hvilke verdier av $x$ er ikke tillatt?
\begin{tasks}(3)
  \task $\dfrac{5}{x-3}$
  \task $\dfrac{x}{x^2-4}$
  \task $\dfrac{x+1}{x(x+2)}$
\end{tasks}

\nivaahode{2}

\oppg{1.4}{2}{\uh}
Forkort. Skriv vilkår.
\begin{tasks}(2)
  \task $\dfrac{x^2-6x+9}{x^2-9}$
  \task $\dfrac{2x^2-8}{x^2+4x+4}$
  \task $\dfrac{x^2+6x+9}{3x+9}$
  \task $\dfrac{5-x}{x-5}$
\end{tasks}

\oppg{1.5}{2}{\uh}
Regn ut. Skriv vilkår.
\begin{tasks}(2)
  \task $\dfrac{x^2-1}{2x} \cdot \dfrac{4x}{x+1}$
  \task $\dfrac{x^2-9}{x+1} : \dfrac{x+3}{2x+2}$
  \task $\dfrac{x^2+2x}{x^2-4} \cdot \dfrac{x-2}{x^2}$
\end{tasks}

\oppg{1.6}{2}{\uh}
Finn \textbf{den første feilen}. Forkort riktig.
\[ \frac{3x^2-12}{x^2+2x} \ \longrightarrow\ \frac{3(x-2)(x+2)}{x(x+2)} \ \longrightarrow\ \frac{3(x-2)}{x} \ \longrightarrow\ \frac{3x-6}{x} \ \longrightarrow\ -3 \]

\nivaahode{3}

\oppg{1.7}{3}{\uh}
Er påstanden sann for alle tillatte verdier av $x$ (alltid), for noen av dem (noen ganger) eller for ingen (aldri)? Begrunn svaret.
\begin{tasks}(1)
  \task $\dfrac{x^2-1}{x-1} = x+1$
  \task $\dfrac{x+2}{x} = 2$
  \task $\dfrac{2x+1}{x} = 2$
\end{tasks}

\oppg{1.8}{3}{\uh}
Lag selv et rasjonalt uttrykk som kan forkortes til $\dfrac{x+1}{x-2}$, og som ikke er definert for $x = 3$.
Forklar hvordan du fant det.

\oppg{1.9}{3}{\uh}
Anna forkorter
\[ \frac{x(x-2)}{(x+2)(x-2)} = \frac{x}{x+2} \]
og skriver vilkåret $x \neq -2$.
Forklar hvorfor hun også må skrive $x \neq 2$, selv om $\dfrac{x}{x+2}$ er definert for $x = 2$.

\begin{sjekk}
\begin{tasks}(2)
  \task Forkort $\dfrac{2x+4}{x^2-4}$. Skriv vilkår.
  \task Hvilken verdi av $x$ er ikke tillatt i $\dfrac{x+1}{x-5}$?
  \task! Forklar med et tall hvorfor $\dfrac{x+4}{x+1}$ ikke er lik $4$.
\end{tasks}
\end{sjekk}

% ---------------------------------------------------------------------

\del{2}{Addere og subtrahere}

\begin{ordliste}
fellesnevner & En nevner som alle brøkene kan skrives med. & \mittsprak\\
minste fellesnevner & Den enkleste fellesnevneren. Den har hver faktor så mange ganger som den står i én av nevnerne. & \mittsprak\\
utvide & Gange teller og nevner med den samme faktoren. & \mittsprak\\
trekke sammen & Skrive flere ledd som én brøk eller ett ledd. & \mittsprak\\
\end{ordliste}

\begin{mustkunne}
\textbf{Samme metode som for tall.} Brøker kan bare legges sammen når nevnerne er like.
\[
\frac{1}{6} + \frac{1}{4} = \frac{1\cdot 2}{6\cdot 2} + \frac{1\cdot 3}{4\cdot 3} = \frac{2+3}{12} = \frac{5}{12}
\qquad\quad
\frac{1}{x} + \frac{1}{x+1} = \frac{1\cdot(x+1)}{x(x+1)} + \frac{1\cdot x}{(x+1)x} = \frac{2x+1}{x(x+1)}
\]
\textbf{Metode:} 1) Faktoriser nevnerne. 2) Finn minste fellesnevner. 3) Utvid hver brøk. 4) Trekk sammen. 5) Forkort hvis mulig.\\
\textbf{Minus foran en brøk:} Sett parentes rundt hele telleren. $\ \dfrac{2}{x} - \dfrac{x-1}{x} = \dfrac{2-(x-1)}{x} = \dfrac{3-x}{x}$
\end{mustkunne}

\begin{eksempel}[1 — to ulike nevnere]
Skriv $\dfrac{2}{x} + \dfrac{3}{x+1}$ som én brøk.

\begin{stegtabell}
Fellesnevner: $x(x+1)$ & Nevnerne har ingen felles faktor. Vilkår: $x \neq 0$, $x \neq -1$.\\
$= \dfrac{2(x+1)}{x(x+1)} + \dfrac{3x}{x(x+1)}$ & Utvid hver brøk.\\
$= \dfrac{2x+2+3x}{x(x+1)} = \dfrac{5x+2}{x(x+1)}$ & Trekk sammen. Ingen felles faktorer å forkorte.\\
Kontroll med $x = 1$: & $\frac{2}{1}+\frac{3}{2} = \frac{7}{2}$ og $\frac{7}{1\cdot 2} = \frac{7}{2}$. Det stemmer.\\
\end{stegtabell}
\end{eksempel}

\begin{eksempel}[2 — minus foran brøk]
Skriv $\dfrac{x}{x^2-4} - \dfrac{1}{x+2}$ så enkelt som mulig.

\begin{stegtabell}
$x^2-4 = (x-2)(x+2)$ & Faktoriser. Vilkår: $x \neq 2$ og $x \neq -2$.\\
Fellesnevner: $(x-2)(x+2)$ & Den andre brøken mangler faktoren $x-2$.\\
$= \dfrac{x}{(x-2)(x+2)} - \dfrac{x-2}{(x-2)(x+2)}$ & Utvid den andre brøken.\\
$= \dfrac{x-(x-2)}{(x-2)(x+2)}$ & Parentes rundt telleren fordi det er minus foran.\\
$= \dfrac{2}{(x-2)(x+2)}$ & $x - x + 2 = 2$.\\
Kontroll med $x = 3$: & $\frac{3}{5} - \frac{1}{5} = \frac{2}{5}$ og $\frac{2}{1\cdot 5} = \frac{2}{5}$. Det stemmer.\\
\end{stegtabell}
\end{eksempel}

\begin{vanligfeil}
\begin{itemize}[leftmargin=1.2em,itemsep=1pt,topsep=0pt]
  \item $\frac{a}{b} + \frac{c}{d} = \frac{a+c}{b+d}$ er feil. Prøv: $\frac{1}{2}+\frac{1}{3} = \frac{5}{6}$, ikke $\frac{2}{5}$.
  \item $\frac{2}{x} - \frac{x-1}{x} = \frac{2-x-1}{x}$ er feil. Minuset gjelder hele telleren: $\frac{2-x+1}{x}$.
  \item $\frac1x+\frac1{x^2}$ med fellesnevner $x^3$ gir $\frac{x^2+x}{x^3}$, som må forkortes til $\frac{x+1}{x^2}$. Minste fellesnevner er $x^2$.
  \item Å glemme å utvide telleren. Når nevneren ganges med $(x+1)$, må telleren også det.
\end{itemize}
\end{vanligfeil}

\nivaahode{1}

\oppg{2.1}{1}{\uh}
Regn ut. Finn fellesnevner først.
\begin{tasks}(3)
  \task $\dfrac{1}{3}+\dfrac{1}{4}$
  \task $\dfrac{5}{6}-\dfrac{3}{8}$
  \task $\dfrac{2}{5}+\dfrac{3}{10}$
\end{tasks}

\oppg{2.2}{1}{\uh}
Trekk sammen.
\begin{tasks}(2)
  \task $\dfrac{3}{x}+\dfrac{5}{x}$
  \task $\dfrac{2x+1}{x+3}+\dfrac{x+5}{x+3}$
  \task $\dfrac{5}{x}-\dfrac{2}{x}$
  \task $\dfrac{x+4}{x-1}-\dfrac{x+1}{x-1}$
\end{tasks}

\oppg{2.3}{1}{\uh}
Finn minste fellesnevner.
\begin{tasks}(3)
  \task $\dfrac{1}{x}$ og $\dfrac{1}{x+2}$
  \task $\dfrac{3}{x^2-4}$ og $\dfrac{1}{x-2}$
  \task $\dfrac{1}{2x}$ og $\dfrac{1}{3x^2}$
\end{tasks}

\nivaahode{2}

\oppg{2.4}{2}{\uh}
Skriv som én brøk. Forkort hvis det går.
\begin{tasks}(2)
  \task $\dfrac{1}{x}+\dfrac{2}{x+3}$
  \task $\dfrac{3}{x-1}-\dfrac{2}{x}$
  \task $\dfrac{1}{4x^2}+\dfrac{1}{6x}$
  \task $\dfrac{x}{x+1}+\dfrac{1}{x-1}$
\end{tasks}

\oppg{2.5}{2}{\uh}
Faktoriser nevnerne først. Skriv som én brøk, forkort og skriv vilkår.
\begin{tasks}(2)
  \task $\dfrac{1}{x-3}-\dfrac{6}{x^2-9}$
  \task $\dfrac{2x}{x^2-1}+\dfrac{1}{x+1}$
  \task $\dfrac{1}{x}-\dfrac{1}{x^2+x}$
\end{tasks}

\oppg{2.6}{2}{\uh}
Finn \textbf{den første feilen}. Regn riktig.
\[ \frac{3}{x-1}-\frac{x+1}{x^2-1} \ \longrightarrow\ \frac{3(x+1)}{(x-1)(x+1)}-\frac{x+1}{(x-1)(x+1)} \ \longrightarrow\ \frac{3x+3-x+1}{(x-1)(x+1)} \ \longrightarrow\ \frac{2x+4}{(x-1)(x+1)} \]

\oppg{2.7}{2}{\mh}
Anna kjører $s$ km til hytta med $60$ km/h og samme vei tilbake med $40$ km/h.
\begin{tasks}(1)
  \task Forklar hvorfor turen tar $\dfrac{s}{60}+\dfrac{s}{40}$ timer. Skriv summen som én brøk.
  \task Gjennomsnittsfart er total strekning delt på total tid. Finn Annas gjennomsnittsfart.
\end{tasks}

\nivaahode{3}

\oppg{2.8}{3}{\uh}
Den ene likningen har ingen løsning. Den andre stemmer for alle tillatte verdier.
Finn ut hvilken som er hvilken, og begrunn.
\begin{tasks}(2)
  \task $\dfrac{x}{x-2} = \dfrac{2}{x-2}+3$
  \task $\dfrac{2}{x-1}+1 = \dfrac{x+1}{x-1}$
\end{tasks}

\oppg{2.9}{3}{\uh}
Er likningen
\[ \frac{1}{a}+\frac{1}{b} = \frac{2}{a+b} \]
sann for noen tillatte verdier av $a$ og $b$? Skriv begge sider som én brøk og begrunn svaret.

\oppg{2.10}{3}{\uh}
Anna kjører strekningen $s$ med fart $v_1$ og samme vei tilbake med fart $v_2$.
Vis at gjennomsnittsfarten alltid er $\dfrac{2v_1v_2}{v_1+v_2}$.

\begin{sjekk}
\begin{tasks}(2)
  \task Skriv $\dfrac{1}{x}+\dfrac{1}{y}$ som én brøk.
  \task Regn ut $\dfrac{2}{x}-\dfrac{3}{x+2}$.
  \task! Hvorfor er $\dfrac{1}{4}+\dfrac{1}{4}$ ikke lik $\dfrac{2}{8}$? Forklar.
\end{tasks}
\end{sjekk}
